\section*{Abstract}

In the context of English as a foreign language (EFL) teaching, classifying texts according to their difficulty level and adapting materials across levels are two central needs for both teachers and learners. This work, carried out within the TaskGen project, investigates and evaluates approaches based on open-source large language models (LLMs) to address both tasks according to the levels of the Common European Framework of Reference for Languages (CEFR).

For the classification task, several open-access LLMs (from the LLaMA, Mistral, and Gemma families) and proprietary models (ChatGPT, DeepSeek) were compared under 24 prompting strategies, including 0-shot, 1-shot, and 2-shot approaches, as well as prompts enriched with domain knowledge such as CEFR descriptors and level-specific vocabulary lists. In addition, the Meta-Llama-3.1-8B-Instruct model was fine-tuned for this task, and a novel metric, Approximate Accuracy, was proposed to capture the ordinal nature of the CEFR scale. Results show that the fine-tuned model consistently outperforms prompting-only strategies and commercial tools such as Text Inspector, reaching a performance comparable to the TSAR reference classifier. Nevertheless, a systematic difficulty remains in distinguishing adjacent levels, particularly at the B1--B2 and B2--C1 boundaries.

For the adaptation task, framed within the TSAR 2025 shared task, the TaskGen system was developed based exclusively on Llama-3.1-8B-Instruct, evaluating three strategies (runs): an iterative refinement focused on content preservation, a refinement step aimed at removing meta-information, and a heuristic voting ensemble over ten prompts of varying complexity. This last strategy achieved the best overall performance, ranking TaskGen 6th among all participating teams, and stood out particularly in meaning preservation (as measured by MeaningBERT) with respect to both the source text and expert-written reference adaptations.

Overall, this work shows that small-scale open-source LLMs (8B parameters), combined with carefully designed prompting strategies and explicit linguistic resources, represent a viable and competitive alternative to larger proprietary models---an outcome particularly relevant for research contexts with limited computational resources, such as the Argentine university setting in which this work was developed.